\section{Affine and metric coordinate geometry}
\label{sec:affine-metric}

This chapter treats points and vectors in dimensions two and three. Unless
explicitly stated otherwise their coordinates, the coefficients of equations,
and construction parameters are rational. Thus every equality test and every
sign case below is a test of integers after multiplication by positive
denominators. Lengths and the square roots occurring in intersection formulas
are interpreted by the fixed primitive recursive real names constructed in the
foundational chapter. The assertions extend to fixed families of such names
when the required denominator bounds and sign certificates are supplied;
they do not assert a decision procedure for equality of arbitrary real names.

Write
\[
 \langle u,v\rangle=\sum_{i=1}^{d}u_iv_i,\qquad
 S(u)=\langle u,u\rangle,\qquad |u|=\sqrt{S(u)},\qquad d\in\{2,3\}.
\]
For planar vectors set
\[
 J(u_1,u_2)=(-u_2,u_1),\qquad [u,v]=u_1v_2-u_2v_1.
\]
A nonzero rational vector has strictly positive squared length: one of its
finitely many rational coordinates is nonzero, its square is positive, and
the other squares are nonnegative. Every reciprocal of a squared length
used below is therefore guarded by a decidable, explicitly stated condition.
In particular, no argument obtains a reciprocal from a mere double negation
of nonzeroness.

\subsection{Incidence and affine constructions}

\begin{theorem}[Equivalent equations of a line, and the line through two points]
\label{geo:line}
Let $P\in\Q^2$ and $u\in\Q^2$ with $u\ne0$. Then
\[
 X=P+tu\quad\Longleftrightarrow\quad [u,X-P]=0
 \quad\Longleftrightarrow\quad \langle Ju,X-P\rangle=0.
\]
Here the existential parameter $t$ is uniquely determined and rational for
rational $X$. Consequently two distinct rational points determine a unique
line, and every equation $ax+by=c$ with $(a,b)\ne(0,0)$ describes a line.
\end{theorem}
\begin{proof}
Substitution proves both implications from the parametric equation. Conversely,
if $u_1\ne0$, set $t=(X_1-P_1)/u_1$. The determinant equation says
$u_1(X_2-P_2)=u_2(X_1-P_1)$, so $X_2-P_2=tu_2$.
If $u_1=0$, then $u_2\ne0$, the determinant equation gives $X_1=P_1$,
and $t=(X_2-P_2)/u_2$ works. The same nonzero coordinate proves uniqueness.
For distinct $P,Q$, take $u=Q-P$. A line containing $P$ and $Q$ has a
direction proportional to $u$, so its point set agrees with this one.
Finally, for $n=(a,b)\ne0$, the rational point
$P=(c/S(n))n$ satisfies $\langle n,P\rangle=c$ and the nonzero direction
$Jn$ is perpendicular to $n$. The equivalence already proved gives the
required line. All branches and quotients have displayed rational guards.
\end{proof}

\begin{theorem}[Slope, intercepts, parallelism, and perpendicularity]
\label{geo:slope}
For $(a,b)\ne(0,0)$, the line $ax+by=c$ has slope-intercept equation
$y=mx+k$ with $m=-a/b$ and $k=c/b$ when $b\ne0$; when $b=0$ it is the
vertical line $x=c/a$. Its intersections with the axes are obtained by
setting the other coordinate equal to zero, with the zero-coefficient case
giving either the whole axis or no point. Two nonvertical lines with slopes
$m,m'$ have parallel directions exactly when $m=m'$, and perpendicular
directions exactly when $mm'=-1$.
\end{theorem}
\begin{proof}
Solving the equation for $y$, or for $x$ if $b=0$, proves the first assertion.
For the $x$-axis the remaining equation is $ax=c$: if $a\ne0$ the unique
intercept is $(c/a,0)$; if $a=0$ its truth is the decidable test $c=0$.
The $y$-axis is identical with the roles reversed. Directions of the two
nonvertical lines can be chosen as $(1,m)$ and $(1,m')$. Their determinant
is $m'-m$ and their dot product is $1+mm'$, proving the last assertions.
The treatment of vertical lines explains why a slope alone cannot represent
every line and why the perpendicular-slope formula requires its stated domain.
\end{proof}

\begin{theorem}[Complete intersection classification of two rational lines]
\label{geo:two-lines}
Consider $ax+by=c$ and $a'x+b'y=c'$, with both normal vectors nonzero, and
put $\Delta=ab'-ba'$. If $\Delta\ne0$, their unique intersection is
\[
 x=\frac{cb'-bc'}{\Delta},\qquad
 y=\frac{ac'-ca'}{\Delta}.
\]
If $\Delta=0$, the lines coincide precisely when
$ac'=a'c$ and $bc'=b'c$; otherwise they are disjoint.
\end{theorem}
\begin{proof}
Multiplying the equations respectively by $b',b$ and subtracting gives
$\Delta x=cb'-bc'$. The analogous elimination gives
$\Delta y=ac'-ca'$. This proves necessity and uniqueness when the determinant
is nonzero, and direct substitution proves existence.
If the determinant vanishes, choose a nonzero coordinate of $(a,b)$.
For example, when $a\ne0$, the determinant equation gives
$(a',b')=\lambda(a,b)$ for $\lambda=a'/a$.
Because the second normal is nonzero, $\lambda\ne0$.
The two equations are then consistent exactly when $c'=\lambda c$,
equivalently $ac'=a'c$. The additional displayed equality follows from it.
The branch $a=0$ uses $b\ne0$ and is identical. An inconsistent proportional
pair cannot have a common point because substitution would force the failed
constant equality. This is a finite rational case split, including coincidence.
\end{proof}

\begin{theorem}[Determinant, collinearity, and triangle area]
\label{geo:area}
For $A,B,C\in\Q^2$, define the signed area
\[
 \mathcal A(A,B,C)=\tfrac12[B-A,C-A].
\]
It is translation invariant, changes sign on interchanging two vertices,
and is zero exactly when the three points are collinear, allowing coincident
points. Its absolute value is one half of base length times perpendicular
height. For every $P$ one has
\[
 \mathcal A(A,B,C)=\mathcal A(P,B,C)+\mathcal A(A,P,C)+\mathcal A(A,B,P).
\]
\end{theorem}
\begin{proof}
Translation cancels in both differences; exchanging $B,C$ reverses the
determinant. Exchanging $A,B$ gives
$[A-B,C-B]=[-(B-A),(C-A)-(B-A)]=-[B-A,C-A]$.
If $A\ne B$, the zero determinant is exactly the line criterion of
Theorem~\ref{geo:line}. If $A=B$, any line through $A,C$ works, and if all
three coincide the horizontal line through them works.
For nonzero base $u=B-A$, decompose $v=C-A$ as
\[
 v=\frac{\langle u,v\rangle}{S(u)}u+
       \frac{[u,v]}{S(u)}Ju.
\]
The identity follows coordinatewise; the second summand is perpendicular
to the base and has length $|[u,v]|/|u|$. Hence half base times height is
$|[u,v]|/2$. A zero base gives zero area directly.
Finally, expand every determinant in the last displayed equation using
$[X-Y,Z-Y]=[X,Z]-[X,Y]-[Y,Z]$; all terms containing $P$ cancel.
This also verifies finite additivity of oriented triangle area without a
topological argument about the inside of a curve.
\end{proof}

\begin{theorem}[Section formula, midpoint, and segment inequalities]
\label{geo:section}
Let $P,Q\in\Q^d$, $P\ne Q$, and $t\in\Q$. The point
$R=(1-t)P+tQ$ is on their line and is on the closed segment precisely when
$0\le t\le1$. In that case
$|R-P|=t|Q-P|$ and $|Q-R|=(1-t)|Q-P|$.
For positive rationals $m,n$, the unique point on the segment with
$|R-P|:|Q-R|=m:n$ is
\[
 R=\frac{nP+mQ}{m+n}.
\]
The midpoint $M=(P+Q)/2$ satisfies, for every $X$,
\[
 S(X-P)+S(X-Q)=2S(X-M)+\tfrac12 S(P-Q).
\]
\end{theorem}
\begin{proof}
The segment is defined parametrically by $0\le t\le1$, and the uniqueness
of the line parameter makes this definition unambiguous. The differences
are $t(Q-P)$ and $(1-t)(Q-P)$, whose lengths have the displayed values by
nonnegativity of $t,1-t$. The ratio equation reduces, after cancelling the
positive base length, to $nt=m(1-t)$, whose unique solution is
$t=m/(m+n)$. For the last identity set $a=X-M$ and $b=(Q-P)/2$.
The left side is $S(a+b)+S(a-b)=2S(a)+2S(b)$, because the two cross terms
cancel. This proves the midpoint identity by a finite polynomial calculation.
\end{proof}

\begin{theorem}[Barycentric coordinates and concurrence of medians]
\label{geo:centroid}
If $A,B,C\in\Q^2$ are noncollinear, every rational point $X$ has unique
rational barycentric coordinates $\alpha,\beta,\gamma$ with
\[
 \alpha+\beta+\gamma=1,\qquad X=\alpha A+\beta B+\gamma C.
\]
They are
\[
 \alpha=\frac{\mathcal A(X,B,C)}{\mathcal A(A,B,C)},\quad
 \beta=\frac{\mathcal A(A,X,C)}{\mathcal A(A,B,C)},\quad
 \gamma=\frac{\mathcal A(A,B,X)}{\mathcal A(A,B,C)}.
\]
The three medians meet uniquely at $G=(A+B+C)/3$, which divides each median
in the ratio $2:1$ measured from its vertex.
\end{theorem}
\begin{proof}
Write $u=B-A$, $v=C-A$. The system $X-A=\beta u+\gamma v$ has a unique
solution because $[u,v]\ne0$. Taking determinants with $v$ and with $u$
gives $\beta=[X-A,v]/[u,v]$ and $\gamma=[u,X-A]/[u,v]$.
Set $\alpha=1-\beta-\gamma$; the area-addition identity proves the formula
for $\alpha$. If $M=(B+C)/2$, then
$G=A+\frac23(M-A)$, so the section formula proves the claimed ratio and
incidence. A point on the median from $A$ has $\beta=\gamma$; a point on
the median from $B$ has $\alpha=\gamma$. Their unique common solution is
$\alpha=\beta=\gamma=1/3$. Thus concurrency and uniqueness are both explicit.
\end{proof}

\begin{theorem}[Affine changes of coordinates]
\label{geo:affine}
Let $T(X)=MX+b$ on $\Q^2$. It preserves affine combinations whose coefficients
sum to one. If $\det M\ne0$, it is bijective, maps lines to lines, preserves
the parameter of every section of a line, and obeys
\[
 [Mu,Mv]=(\det M)[u,v],\qquad
 \mathcal A(TA,TB,TC)=(\det M)\mathcal A(A,B,C).
\]
For a locus $F(X)=0$, its image is described by
$F(M^{-1}(Y-b))=0$.
\end{theorem}
\begin{proof}
Linearity gives
$T(\sum_i\lambda_iX_i)=\sum_i\lambda_iT(X_i)$ when $\sum_i\lambda_i=1$.
For $M=\left(\begin{smallmatrix}a&b_0\\c&d\end{smallmatrix}\right)$,
expansion gives
\[
 (au_1+b_0u_2)(cv_1+dv_2)-(cu_1+du_2)(av_1+b_0v_2)
 =(ad-b_0c)(u_1v_2-u_2v_1).
\]
The inverse matrix is
$(ad-b_0c)^{-1}\left(\begin{smallmatrix}d&-b_0\\-c&a\end{smallmatrix}\right)$,
verified by multiplication. Thus a nonzero direction stays nonzero, and
$T(P+tu)=T(P)+tMu$ proves the assertions about lines and sections.
The area formula follows from the determinant identity. Finally
$Y=T(X)$ is equivalent to $X=M^{-1}(Y-b)$, proving both directions of
the locus transformation. The inverse substitution is essential: substituting
the forward map into $F$ would describe its inverse image.
\end{proof}

\subsection{Inner products, distances, projections, and isometries}

\begin{theorem}[Lagrange's sum-of-squares identity]
\label{geo:lagrange}
In dimensions two and three,
\[
 S(u)S(v)-\langle u,v\rangle^2
 =\sum_{1\le i<j\le d}(u_iv_j-u_jv_i)^2.
\]
In the plane the right side is $[u,v]^2$. It vanishes precisely when $u,v$
are linearly dependent, including the case where one vector is zero.
\end{theorem}
\begin{proof}
Expand the product as
$\sum_i u_i^2v_i^2+\sum_{i<j}(u_i^2v_j^2+u_j^2v_i^2)$.
The square of the dot product is
$\sum_i u_i^2v_i^2+2\sum_{i<j}u_iu_jv_iv_j$.
Subtracting leaves exactly the displayed squares. If they all vanish and
$u\ne0$, choose $k$ with $u_k\ne0$. Each equation
$u_kv_j-u_jv_k=0$ gives $v_j=(v_k/u_k)u_j$, proving dependence.
If $u=0$, dependence is immediate. Conversely proportional vectors make
every minor vanish. A finite sum of nonnegative rational squares is zero
only if each square is zero, a fact proved by rational order and the finite
number of summands.
\end{proof}

\begin{theorem}[Cauchy--Schwarz and its equality condition]
\label{geo:cauchy}
For $u,v\in\Q^d$,
\[
 |\langle u,v\rangle|\le |u||v|.
\]
Equality holds exactly when the vectors are linearly dependent. For nonzero
vectors the normalized dot product lies between $-1$ and $1$.
\end{theorem}
\begin{proof}
The sum of squares in Theorem~\ref{geo:lagrange} is nonnegative, so
$\langle u,v\rangle^2\le S(u)S(v)$. Both sides of the desired inequality
are nonnegative; comparison of their squares proves it, using monotonicity
of the nonnegative square root. Equality is equivalent to the vanishing
of the difference of those squares, hence to linear dependence by the same
theorem. If both squared lengths are positive, their product has a positive
square root, so division gives the final assertion. For rational inputs this
positive denominator has an explicit rational lower bound, obtainable from
the input numerators and denominators or the square-root construction.
\end{proof}

\begin{theorem}[Pythagoras, parallelogram identity, and polarization]
\label{geo:pythagoras}
The identities
\[
 S(u+v)=S(u)+2\langle u,v\rangle+S(v),\qquad
 S(u+v)+S(u-v)=2S(u)+2S(v)
\]
hold in dimensions two and three. In particular,
$S(u+v)=S(u)+S(v)$ exactly when $\langle u,v\rangle=0$.
Moreover
\[
 \langle u,v\rangle=\tfrac14\bigl(S(u+v)-S(u-v)\bigr).
\]
\end{theorem}
\begin{proof}
For each coordinate,
$(u_i+v_i)^2=u_i^2+2u_iv_i+v_i^2$.
Summation proves the first identity. Replacing $v$ by $-v$ and adding cancels
the dot products; subtracting instead yields four times the dot product.
The Pythagorean equivalence follows by subtracting $S(u)+S(v)$ and using
the nonzero rational constant $2$. Thus the usual geometric right-angle
criterion, defined by a zero inner product, is equivalent to the squared
length identity; no angle-selection operation is needed.
\end{proof}

\begin{theorem}[The Euclidean distance is a metric]
\label{geo:metric}
Define $d(P,Q)=\sqrt{\sum_i(P_i-Q_i)^2}$. It is nonnegative and symmetric,
equals zero exactly when $P=Q$, and satisfies
\[
 d(P,R)\le d(P,Q)+d(Q,R).
\]
For a nonnegative rational $a$, $d(aP,aQ)=a\,d(P,Q)$.
\end{theorem}
\begin{proof}
Nonnegativity is part of the square-root specification. Reversing each
difference leaves its square unchanged, proving symmetry. If the distance
is zero, its square is zero, and each nonnegative coordinate square is zero;
therefore all coordinate differences vanish. The converse is immediate.
For the triangle inequality put $u=P-Q$ and $v=Q-R$. By the preceding
theorems,
\[
 S(u+v)=S(u)+2\langle u,v\rangle+S(v)
 \le S(u)+2|u||v|+S(v)=(|u|+|v|)^2.
\]
Taking nonnegative square roots gives $|u+v|\le|u|+|v|$.
Finally $S(a(P-Q))=a^2S(P-Q)$, and $a|P-Q|$ is the unique nonnegative
square root of that value. Every appeal to square roots is an inequality
or equality of the represented values; no branching on their equality occurs.
\end{proof}

\begin{theorem}[Cosine law and an algebraic encoding of planar angles]
\label{geo:cosine}
For a triangle with side vectors $u,v$ at one vertex and nonzero adjacent
lengths $a=|u|$, $b=|v|$, set
$c_\theta=\langle u,v\rangle/(ab)$. Then
\[
 |u-v|^2=a^2+b^2-2ab c_\theta.
\]
In the plane, set also $s_\theta=[u,v]/(ab)$. Then
$c_\theta^2+s_\theta^2=1$. These two numbers encode an oriented angle
without choosing an inverse trigonometric branch. For nonvertical directions
$(1,m),(1,m')$, when $1+mm'\ne0$, the corresponding sine-to-cosine ratio is
\[
 \frac{s_\theta}{c_\theta}=\frac{m'-m}{1+mm'}.
\]
\end{theorem}
\begin{proof}
The cosine law is the expansion of $S(u-v)$ followed by the definition of
$c_\theta$. The planar Lagrange identity gives
$\langle u,v\rangle^2+[u,v]^2=S(u)S(v)=a^2b^2$;
division proves the unit-circle equation. The dot product of the two slope
directions is $1+mm'$ and their determinant is $m'-m$, so their common
normalizing factor cancels in the displayed quotient. The nonzero-denominator
condition excludes a right angle, where this quotient has no finite value.
The normalized pair remains defined there and also records the quadrant
information lost by an unqualified arctangent of a slope.
\end{proof}

\begin{theorem}[Projection onto a line and its minimum-distance property]
\label{geo:projection}
Let $L=\{P+tu:t\in\Q\}$ with $u\ne0$, and let $X\in\Q^d$.
The perpendicular foot is
\[
 H=P+\frac{\langle X-P,u\rangle}{S(u)}u.
\]
It is the unique point of the line minimizing squared distance to $X$, since
for every parameter $t$,
\[
 S\bigl(X-(P+tu)\bigr)
 =S(X-H)+S(u)\left(t-\frac{\langle X-P,u\rangle}{S(u)}\right)^2.
\]
In the plane,
$S(X-H)=[u,X-P]^2/S(u)$; in three dimensions it is
$S(X-P)-\langle X-P,u\rangle^2/S(u)$.
\end{theorem}
\begin{proof}
The chosen coefficient makes $\langle X-H,u\rangle=0$.
Write $X-(P+tu)=(X-H)+(t_0-t)u$, where $t_0$ is that coefficient.
Pythagoras gives the displayed decomposition. Its second term is nonnegative
and vanishes only when $t=t_0$, proving minimum and uniqueness.
Expanding $S((X-P)-t_0u)$ gives
$S(X-P)-\langle X-P,u\rangle^2/S(u)$. The planar Lagrange identity rewrites
this as the determinant quotient. The same minimum statement holds for fixed
PR real line parameters, since its proof is an exact identity plus a square.
\end{proof}

\begin{theorem}[Projection onto a line or plane in normal form]
\label{geo:hyperplane}
Let $n\ne0$ and $\Pi=\{Y:\langle n,Y\rangle=c\}$ in dimension two or three.
For rational $X$, set
\[
 H=X-\frac{\langle n,X\rangle-c}{S(n)}n.
\]
Then $H\in\Pi$ and, for every $Y\in\Pi$,
\[
 S(X-Y)=\frac{(\langle n,X\rangle-c)^2}{S(n)}+S(H-Y).
\]
Thus $H$ is the unique nearest point and
$d(X,\Pi)=|\langle n,X\rangle-c|/|n|$.
\end{theorem}
\begin{proof}
Taking the dot product with $n$ proves $\langle n,H\rangle=c$.
The vector $X-H$ is a scalar multiple of $n$, while
$\langle n,H-Y\rangle=0$. Applying Pythagoras to
$X-Y=(X-H)+(H-Y)$ gives the identity, and
$S(X-H)=(\langle n,X\rangle-c)^2/S(n)$.
The nonnegative second term vanishes exactly when $Y=H$, proving the claims.
This treats point-to-line and point-to-plane distance with one construction.
\end{proof}

\begin{theorem}[Reflection in a line or plane]
\label{geo:reflection}
For the normal equation in Theorem~\ref{geo:hyperplane}, the map
\[
 R(X)=X-2\frac{\langle n,X\rangle-c}{S(n)}n
\]
is an involution fixing exactly $\Pi$, preserves all distances, and has
midpoint $(X+R(X))/2$ equal to the perpendicular foot of $X$.
\end{theorem}
\begin{proof}
Let $s=\langle n,X\rangle-c$. Then
$\langle n,R(X)\rangle-c=-s$, so applying the same formula twice returns
$X$. The equation $R(X)=X$ implies $s=0$ by taking a nonzero coordinate of
$n$; the converse is immediate. The midpoint formula is the projection
formula. For a difference vector $v=X-Y$, its image is
$v-2\langle n,v\rangle n/S(n)$. Expanding its squared length gives
\[
 S(v)-4\frac{\langle n,v\rangle^2}{S(n)}
       +4\frac{\langle n,v\rangle^2}{S(n)}=S(v).
\]
Taking nonnegative square roots proves preservation of distance.
\end{proof}

\begin{theorem}[Orthogonal transformations and similarities]
\label{geo:isometries}
If a rational matrix $M$ satisfies $M^{\mathsf T}M=I$, then $X\mapsto MX+b$
preserves dot products of differences and distances. Its inverse is
$Y\mapsto M^{\mathsf T}(Y-b)$. A rational linear map preserves every squared
length if and only if $M^{\mathsf T}M=I$. In dimension two its determinant
is $1$ or $-1$, and it has one of the forms
\[
 \begin{pmatrix}c&-s\\s&c\end{pmatrix},\qquad
 \begin{pmatrix}c&s\\s&-c\end{pmatrix},\qquad c^2+s^2=1.
\]
Multiplying an orthogonal matrix by a nonzero rational $a$ gives a similarity
with distance factor $|a|$.
\end{theorem}
\begin{proof}
Matrix multiplication gives
$\langle Mu,Mv\rangle=u^{\mathsf T}M^{\mathsf T}Mv=\langle u,v\rangle$.
An orthogonal square matrix has nonzero determinant because
$(\det M)^2=1$; the inverse follows. Conversely, squared-length preservation
on each standard basis vector makes each column have squared length one;
preservation on the sum of each pair makes their dot product zero by
polarization. These finitely many equations are exactly $M^{\mathsf T}M=I$.
In dimension two write the first column as $(c,s)$, of squared length one.
Its orthogonal complement consists of multiples of $(-s,c)$ by the line
criterion. A unit second column is therefore $t(-s,c)$ with $t^2=1$.
The rational factorization $(t-1)(t+1)=0$ gives $t=1$ or $t=-1$, yielding
the two forms. Finally $S(aMv)=a^2S(v)$ proves the similarity factor.
Only the rational parameters are classified; there is no algorithmic
classification of arbitrary real matrix entries by unbounded equality tests.
\end{proof}

\subsection{Three-dimensional incidence}

For $u,v\in\Q^3$, define
\[
 u\mathbin{\times}v=(u_2v_3-u_3v_2,\ u_3v_1-u_1v_3,\ u_1v_2-u_2v_1).
\]

\begin{theorem}[Cross product, triple product, and volume]
\label{geo:cross}
The cross product is perpendicular to both factors and satisfies
\[
 S(u\times v)=S(u)S(v)-\langle u,v\rangle^2,\qquad
 u\times(v\times w)=v\langle u,w\rangle-w\langle u,v\rangle.
\]
The scalar triple product $\langle u\times v,w\rangle$ is the determinant
with columns $u,v,w$; it vanishes precisely when those vectors are dependent.
The absolute volume of their parallelepiped is its absolute value.
\end{theorem}
\begin{proof}
In $\langle u,u\times v\rangle$ the six monomials cancel in pairs, and
the same happens with $v$. The squared-length formula is precisely the
three-dimensional Lagrange identity, with its three minors in the stated
order. For the first coordinate of the vector triple identity, expansion gives
\[
 u_2(v_1w_2-v_2w_1)-u_3(v_3w_1-v_1w_3)
 =v_1(u_2w_2+u_3w_3)-w_1(u_2v_2+u_3v_3).
\]
Adding and subtracting the equal terms $u_1v_1w_1$ supplies the first
coordinate on the right. The other two coordinates follow by cyclic renaming.
Expansion of the scalar triple product is the six-term determinant formula.
If $u\times v\ne0$, choose a nonzero $2\times2$ minor of $(u,v)$.
It solves two coordinates of $w=\alpha u+\beta v$, and the zero triple
product supplies the third; if $u\times v=0$, the first two vectors are
already dependent by Theorem~\ref{geo:lagrange}. Conversely dependence
forces determinant zero by linearity and alternating signs.
When $u\times v\ne0$, the base parallelogram has area $|u\times v|$ by
the base-times-height calculation, and the component of $w$ normal to it
has length $|\langle u\times v,w\rangle|/|u\times v|$.
Their product is the asserted volume. A dependent base gives zero volume.
\end{proof}

\begin{theorem}[Equations of a plane and a plane through three points]
\label{geo:plane}
For $n\in\Q^3\setminus\{0\}$ and $P\in\Q^3$, the plane with normal $n$
through $P$ has equivalent equations
\[
 \langle n,X-P\rangle=0,\qquad
 n_1x+n_2y+n_3z=\langle n,P\rangle.
\]
It has a rational parametrization by two independent direction vectors.
Three noncollinear rational points $A,B,C$ determine the unique plane with
normal $(B-A)\times(C-A)$ through $A$.
\end{theorem}
\begin{proof}
Expansion gives equivalence of the two equations. If $n_3\ne0$, take
directions $u=(1,0,-n_1/n_3)$ and $v=(0,1,-n_2/n_3)$.
Every point in the plane has the unique form $P+su+tv$, because its first
two differences determine $s,t$ and the normal equation forces the third.
If $n_3=0$, use a nonzero coordinate among $n_1,n_2$ and permute this
construction. The two displayed directions are independent.
For the three-point assertion their differences are independent, so their
cross product is nonzero. It is perpendicular to both differences and
therefore defines a plane containing all three. Any other plane containing
the points has normal perpendicular to both independent directions; solving
the same two-coordinate system shows that normal is proportional to their
cross product. Consequently its equation describes the same plane.
\end{proof}

\begin{theorem}[Line--plane intersection, including all degenerate cases]
\label{geo:line-plane}
Let $X=P+tu$ be a line with $u\ne0$, and let
$\langle n,X\rangle=c$ be a plane with $n\ne0$.
Put $a=\langle n,u\rangle$ and $b=c-\langle n,P\rangle$.
If $a\ne0$, their unique intersection occurs at $t=b/a$.
If $a=0=b$, the whole line lies in the plane; if $a=0$ and $b\ne0$,
there is no intersection.
\end{theorem}
\begin{proof}
Substituting the line equation gives exactly the scalar equation $at=b$.
The three rational cases give its complete solution set. In the first case
the nonzero direction ensures a unique point corresponds to the unique
parameter. Thus this classification is a bounded amount of rational arithmetic
with no appeal to a general equation-solving principle.
\end{proof}

\begin{theorem}[Intersection of two planes]
\label{geo:two-planes}
Let the planes be $\langle n,X\rangle=c$ and $\langle m,X\rangle=d$, with
both normals nonzero, and put $k=n\times m$.
If $k\ne0$, their intersection is the line
\[
 X=X_0+tk,\qquad
 X_0=\frac{c(m\times k)+d(k\times n)}{S(k)}.
\]
If $k=0$, their normals are proportional; they either coincide or are
disjoint, according as their constants have the same proportionality factor.
\end{theorem}
\begin{proof}
The triple-product identities give
$\langle n,m\times k\rangle=S(k)$,
$\langle m,k\times n\rangle=S(k)$, and the two crossed dot products
$\langle n,k\times n\rangle$, $\langle m,m\times k\rangle$ equal zero.
Hence $X_0$ satisfies both equations. The direction $k$ is perpendicular
to both normals, so the whole displayed line lies in the intersection.
Conversely, if $X$ is another solution, $X-X_0$ is perpendicular to the
independent vectors $n,m$. Solving using any nonzero minor of $(n,m)$
shows that their common orthogonal complement consists exactly of multiples
of $k$. If $k=0$, Theorem~\ref{geo:lagrange} yields $m=\lambda n$ with
nonzero rational $\lambda$, obtained from a nonzero coordinate of $n$.
The equations then coincide exactly when $d=\lambda c$; otherwise
substitution proves inconsistency.
\end{proof}

\begin{theorem}[Common perpendicular and distance between two spatial lines]
\label{geo:skew}
Let $L=P+tu$ and $L'=Q+sv$ with $u,v\ne0$ and set $w=Q-P$.
If $k=u\times v\ne0$, the minimum squared distance is
\[
 \frac{\langle w,k\rangle^2}{S(k)}.
\]
The unique minimizing parameter pair solves
\[
 \begin{pmatrix}S(u)&-\langle u,v\rangle\\
 -\langle u,v\rangle&S(v)\end{pmatrix}
 \binom{t}{s}=
 \binom{\langle u,w\rangle}{-\langle v,w\rangle}.
\]
The lines intersect exactly when $\langle w,k\rangle=0$; otherwise they
are skew. If $k=0$, their distance is the point-to-line distance from $Q$
to $L$, and they coincide exactly when that squared distance is zero.
\end{theorem}
\begin{proof}
For a minimizing pair the joining vector
$D=(Q+sv)-(P+tu)=w+sv-tu$ must be perpendicular to both directions.
Those two perpendicularity equations are exactly the displayed system.
Its determinant is $S(u)S(v)-\langle u,v\rangle^2=S(k)>0$, so Cramer's
formulas construct a unique rational solution $(t_0,s_0)$. Its joining
vector $D_0$ is a multiple of $k$. Taking its dot product with $k$ gives
$\langle D_0,k\rangle=\langle w,k\rangle$, whence
$D_0=\langle w,k\rangle k/S(k)$ and the distance formula.
For any other parameters,
$D=D_0+(s-s_0)v-(t-t_0)u$, with the last vector perpendicular to $D_0$.
Pythagoras gives $S(D)\ge S(D_0)$, with equality only if the last vector
vanishes. Independence of $u,v$ forces the same pair of parameters.
A zero minimum is an actual intersection because the minimizing pair is
constructed; positive minimum proves disjointness. If $k=0$, write
$v=\lambda u$ with $\lambda\ne0$. The family of joining vectors is then
$w+ru$, and its minimum is the projection construction of
Theorem~\ref{geo:projection}. A zero value means $Q$ is on $L$, which
makes the two proportional-direction lines identical.
\end{proof}

\subsection{Circles and elementary triangle geometry}

\begin{theorem}[Circle equation and completion of the square]
\label{geo:circle}
A circle with rational center $O=(h,k)$ and squared radius $q\ge0$ has
equation $(x-h)^2+(y-k)^2=q$.
The equation $x^2+y^2+ax+by+c=0$ is equivalent to this form with
\[
 O=(-a/2,-b/2),\qquad q=(a^2+b^2)/4-c.
\]
It has no real points when $q<0$, one point when $q=0$, and a nontrivial
circle when $q>0$, with explicit points $O\pm(\sqrt q,0)$.
\end{theorem}
\begin{proof}
Expand $(x+a/2)^2+(y+b/2)^2$ and move its constant term to the other
side. The left side is a sum of two nonnegative squares, so a negative
right side is impossible. A zero right side forces both squares to vanish,
giving exactly the center. For positive $q$ the indicated points satisfy
the equation and are distinct because $\sqrt q>0$. Thus nonemptiness does
not rely on an unbounded search. Rational $q$ need not have a rational square
root; the points are represented by the foundational PR square-root names.
\end{proof}

\begin{theorem}[Line--circle intersections]
\label{geo:line-circle}
Let a circle have rational center $O$ and rational squared radius $q\ge0$,
and let $L=P+tu$ with $u\ne0$. Define
\[
 t_0=\frac{\langle O-P,u\rangle}{S(u)},\qquad H=P+t_0u,\qquad
 K=q-S(H-O).
\]
There are no intersections if $K<0$, exactly one, $H$, if $K=0$, and
exactly two if $K>0$, namely
\[
 H\mathbin{\pm}\sqrt{K/S(u)}\,u.
\]
When $q>0$ and $K=0$, the line is perpendicular to the radius at the
contact point and is the tangent line there.
\end{theorem}
\begin{proof}
The projection formula makes $H-O$ perpendicular to $u$. Every point of
the line is uniquely $H+ru$, so Pythagoras changes its circle equation into
$S(H-O)+r^2S(u)=q$, or $r^2=K/S(u)$.
Its right side is rational and has the three decidable sign cases in the
statement. A negative value is impossible, zero forces $r=0$, and a positive
value has precisely the two roots $\pm\sqrt{K/S(u)}$.
For completeness, write $s=\sqrt{K/S(u)}>0$ and take the positive rational
lower bound $a=\min\{1,K/S(u)\}\le s$. The circle equation
gives $(r-s)(r+s)=0$. Approximate a proposed represented-real parameter $r$
by a rational $r_0$ with $|r-r_0|\le a/4$. If $r_0\ge0$, then
$r+s\ge3a/4$, so multiplication by its guarded reciprocal gives $r=s$.
If $r_0<0$, then $s-r\ge3a/4$, and the same argument gives $r=-s$.
The only test was the sign of the rational approximant; no real sign
trichotomy or general disjunction from a zero product was used.
In the zero case with $q>0$, $H-O\ne0$ and the line has this radius as
a normal, establishing the tangent equation.
\end{proof}

\begin{theorem}[Radical axis and complete circle--circle intersection test]
\label{geo:circle-circle}
Let two circles have rational centers $O_1,O_2$ and squared radii
$q_1,q_2\ge0$. If $O_1\ne O_2$, put
\[
 w=O_2-O_1,\quad D=S(w),\quad
 t=\frac{D+q_1-q_2}{2D},\quad H=O_1+tw,\quad K=q_1-t^2D.
\]
Their common points lie on the radical axis
\[
 \langle X-O_1,w\rangle=(D+q_1-q_2)/2,
\]
and are absent for $K<0$, equal to $H$ alone for $K=0$, and precisely
\[
 H\mathbin{\pm}\sqrt{K/D}\,Jw
\]
for $K>0$. If $O_1=O_2$, the circles coincide when $q_1=q_2$ and otherwise
have no common point.
\end{theorem}
\begin{proof}
Writing $Y=X-O_1$, subtract the equations $S(Y)=q_1$ and
$S(Y-w)=q_2$. Since $S(Y-w)=S(Y)-2\langle Y,w\rangle+D$, subtraction
gives the radical-axis equation. Conversely, that equation together with
the first circle equation implies the second, so no solutions have been
lost. The axis is $H+rJw$, because $\langle H-O_1,w\rangle=tD$ and
$Jw$ is perpendicular to $w$. Its substitution into the first circle gives
$t^2D+r^2D=q_1$, whence the stated three cases by the preceding theorem.
If the centers agree, the equations have the same left side; equality or
inequality of the rational right sides gives the final classification.
The single-point case is called tangency when both radii are positive;
the statement also correctly includes circles of radius zero.
\end{proof}

\begin{theorem}[Circumcenter of a noncollinear rational triangle]
\label{geo:circumcenter}
Let $A,B,C\in\Q^2$ be noncollinear and put $u=B-A$, $v=C-A$.
There is a unique circle through the three vertices. Its rational center is
\[
 O=A+\frac{S(v)Ju-S(u)Jv}{2[u,v]},
\]
and its squared radius is the positive rational $q=S(O-A)$.
The perpendicular bisectors of all three sides pass through $O$.
\end{theorem}
\begin{proof}
For $Y=O-A$, equality of the squared distances to $A,B$ is
$S(Y)=S(Y-u)$, equivalently $2\langle Y,u\rangle=S(u)$.
The analogous equation for $A,C$ is
$2\langle Y,v\rangle=S(v)$.
The two normals are independent. The displayed formula is their unique
solution: $\langle Ju,u\rangle=0$,
$\langle Jv,u\rangle=-[u,v]$, and
$\langle Ju,v\rangle=[u,v]$ verify it directly.
The locus of points equidistant from two distinct points is their
perpendicular bisector, by expanding the same squared-distance equation
around their midpoint. Thus all three bisectors contain $O$.
The radius cannot be zero: that would force $O=A=B=C$, contradicting
noncollinearity. Any other circle through the vertices would have a center
solving the same unique pair of linear equations and the same radius.
\end{proof}

\begin{theorem}[Orthocenter and the Euler-line relation]
\label{geo:orthocenter}
For the triangle of Theorem~\ref{geo:circumcenter}, the three altitudes
meet uniquely at
\[
 H=A+B+C-2O.
\]
With centroid $G=(A+B+C)/3$, one has $H=3G-2O$.
Consequently $O,G,H$ are collinear; if they are not all equal, $G$ divides
the segment from $O$ to $H$ in the ratio $1:2$.
\end{theorem}
\begin{proof}
Translate the origin to $O$, so the three vertex vectors all have squared
length $q$ and the proposed orthocenter is $H=A+B+C$.
Then
\[
 \langle H-A,B-C\rangle
 =\langle B+C,B-C\rangle=S(B)-S(C)=0.
\]
Thus $H$ lies on the altitude through $A$, and cyclic permutation proves
the other two incidences. The normals of the altitudes through $A,B$ are
the nonparallel side directions $B-C,A-C$, so their intersection is unique.
Restoring the original origin gives the displayed formula.
The relation $H=3G-2O$ follows by substitution, equivalently
$G=\frac23O+\frac13H$. The section formula then gives the ratio whenever
the segment has nonzero length. This explicitly includes the exceptional
case where the Euler line has no distinguished direction because its three
named points coincide.
\end{proof}

\begin{theorem}[The nine-point circle]
\label{geo:ninepoint}
The three side midpoints, the three midpoints between a vertex and the
orthocenter, and the three altitude feet of a noncollinear triangle lie on
the circle with center $N=(O+H)/2$ and squared radius $q/4$, where $O,H,q$
are its circumcenter, orthocenter, and squared circumradius.
Repeated named points are allowed.
\end{theorem}
\begin{proof}
Again put $O=0$, so $H=A+B+C$ and $S(A)=S(B)=S(C)=q$.
For the side midpoint $U=(B+C)/2$,
$U-N=-A/2$, so $S(U-N)=q/4$.
For $V=(A+H)/2$, one has $V-N=A/2$, yielding the same equation.
Cyclic permutation treats the six midpoints.
Let $D$ be the foot of the altitude from $A$ onto $BC$, constructed by
Theorem~\ref{geo:projection}. Both $D,U$ lie on $BC$, and both $D,V$
lie on the altitude through $A$. These lines are perpendicular, so
$\langle D-U,D-V\rangle=0$. For any $U,V,D$, writing
$M=(U+V)/2$ and expanding gives
\[
 \langle D-U,D-V\rangle=S(D-M)-\tfrac14 S(U-V).
\]
In this instance $M=N$ and $V-U=A$, so the zero dot product gives
$S(D-N)=q/4$. The other two feet follow by cyclic permutation.
This proof handles obtuse and right triangles without an interior-point
assumption or a special position diagram.
\end{proof}

\begin{theorem}[Power of a point, secants, and equal tangents]
\label{geo:power}
For the circle $S(X-O)=q$ with $q>0$, define the power of $P$ by
$\operatorname{Pow}(P)=S(P-O)-q$.
For a line $P+tu$ with $u\ne0$, suppose its restricted quadratic has two
real roots $t_1,t_2$, counted with multiplicity. Then
\[
 S(u)t_1t_2=\operatorname{Pow}(P).
\]
If $T$ is a contact point and $\langle T-O,P-T\rangle=0$, then
$S(P-T)=\operatorname{Pow}(P)$; hence the lengths of two tangent segments
from the same point are equal. If $D=S(P-O)>q$, two such contact points
are explicitly
\[
 T_\pm=O+\frac qD(P-O)
       \mathbin{\pm}\frac{\sqrt{q(D-q)}}D J(P-O).
\]
\end{theorem}
\begin{proof}
Substitution of the secant into the circle yields
$S(u)t^2+2\langle P-O,u\rangle t+\operatorname{Pow}(P)=0$.
When there are two roots, including a repeated root, comparison with
$S(u)(t-t_1)(t-t_2)$ gives the product identity; it also follows directly
from the quadratic root formulas in Theorem~\ref{geo:line-circle}.
The product is directed: replacing it by an unsigned length product would
lose the sign of the power for a point inside the circle.
For a contact point, Pythagoras applied to
$P-O=(P-T)+(T-O)$ gives $S(P-O)=S(P-T)+q$.
For the final construction put $Y=P-O$. Orthogonality of $Y,JY$ gives
\[
 S(T_\pm-O)=q^2/D+q(D-q)/D=q,
 \qquad \langle Y,T_\pm-O\rangle=q.
\]
Subtracting $S(T_\pm-O)=q$ from the last equality proves
$\langle T_\pm-O,P-T_\pm\rangle=0$.
Because $q(D-q)>0$ and $Y\ne0$, the points are distinct. Their tangent
lines have the respective nonzero radius vectors as normals, and the
line--circle theorem proves that each meets the circle only at its contact point.
\end{proof}

\begin{theorem}[Shoelace formula and origin-independent polygon area]
\label{geo:shoelace}
For a finite cyclic list of rational planar points
$P_0,\ldots,P_{m-1}$, with $m\ge1$ and $P_m=P_0$, define its signed fan area from $O$ by
$\sum_{i<m}\mathcal A(O,P_i,P_{i+1})$. It is independent of $O$ and equals
\[
 \frac12\sum_{i<m}(x_iy_{i+1}-y_ix_{i+1}).
\]
If the polygon comes with an oriented triangulation whose internal edges
occur in opposite directions, this number is the sum of the oriented areas
of its triangles. In particular, for a convex polygon in cyclic positive
order it is its ordinary area.
\end{theorem}
\begin{proof}
Expand
$[P_i-O,P_{i+1}-O]=[P_i,P_{i+1}]+[O,P_i]-[O,P_{i+1}]$.
Summing around the cycle cancels the last two terms, including the final
term because $P_m=P_0$. This proves origin independence and the coordinate
formula. A triangle has twice its signed area equal to the sum of its three
directed-edge determinants, by Theorem~\ref{geo:area}. Summing these
identities over the supplied triangulation cancels every paired internal
edge, leaving the boundary sum. For a convex polygon, the fan from $P_0$
has triangles with nonnegative orientation and covers it without overlap
of interiors; its area is therefore the sum of their base-times-height
areas. Alternatively that finite fan can be taken as the area certificate.
The proof of cancellation for a variable list length is induction on the
length of a coded list with a rational equality invariant, hence quantifier-free
induction after coding. No decision procedure for arbitrary Jordan regions
is asserted.
\end{proof}

\begin{theorem}[Intersection of a sphere with a plane]
\label{geo:sphere-plane}
Let $S(X-O)=q$ with $q\ge0$ be a sphere with rational center in $\Q^3$,
and let $\langle n,X\rangle=c$, $n\ne0$, be a rational plane. Set
\[
 H=O+\frac{c-\langle n,O\rangle}{S(n)}n,\qquad
 K=q-\frac{(c-\langle n,O\rangle)^2}{S(n)}.
\]
The intersection is empty for $K<0$, is the point $H$ for $K=0$, and is
a circle in that plane with center $H$ and squared radius $K$ for $K>0$.
An orthonormal coordinate pair for that circle is constructible using
rational arithmetic and positive square roots.
\end{theorem}
\begin{proof}
The point $H$ is the projection of $O$ onto the plane. Every point of the
plane is $H+v$ with $\langle n,v\rangle=0$; since $H-O$ is a multiple
of $n$, Pythagoras gives
$S(H+v-O)=S(H-O)+S(v)$. Thus the sphere equation is exactly $S(v)=K$,
which proves the three cases except for the promised explicit coordinates.
Choose a standard basis vector $e_i$ with $E=n\times e_i\ne0$.
Such an index exists: a nonzero vector cannot be parallel to all three
basis vectors, and the three rational tests are finite. Put $F=n\times E$.
Then $E,F$ are perpendicular to $n$ and to each other, and
$S(F)=S(n)S(E)>0$ by the cross-product identity. Therefore
\[
 e=E/\sqrt{S(E)},\qquad f=F/\sqrt{S(n)S(E)}
\]
are a constructed orthonormal basis of the plane's direction space.
The circle is exactly $H+ae+bf$ with $a^2+b^2=K$.
All reciprocal square roots have positive rational inputs and explicit
apartness bounds. In particular, when $K>0$ the points $H\pm\sqrt K\,e$
witness nonemptiness constructively.
\end{proof}

\subsection{Arithmetic audit of the chapter}

The thirty statements above use three kinds of proof steps. First, fixed
coordinate identities are finite polynomial identities over rational codes;
clearing their positive denominators reduces them to integer identities, and
splitting signed integers into natural-number codes gives quantifier-free
arithmetic formulas. Second, every rational choice is among a fixed finite
list: a nonzero coordinate, a nonzero minor, a determinant sign, or one of the
three signs of a rational squared-radius difference. Its branch condition is
decidable and its branch result is a primitive recursive expression. Third,
lengths and intersection coordinates use the fixed square-root and arithmetic
constructions already proved in the foundation. Their equality and order
conclusions are the corresponding finite-precision schemas with explicit
error budgets. Polynomial equalities never require deciding equality of their
represented real values.

The only variable-length construction introduced here is the polygon sum.
Its recurrence stores a rational accumulator and a bounded index; the proof
invariant is the displayed telescoping identity. The dimensions in every
matrix calculation remain two or three. Thus the argument does not invoke
unbounded search, an arbitrary-dimensional choice of bases, unrestricted
induction on formulas with quantifiers, or a total evaluator for all primitive
recursive programs. These observations specify the actual arithmetic
translation needed for the Lean development, rather than inferring an
\ipra\ proof merely from the fact that a statement is true over real numbers.
